\documentclass[11pt,a4paper]{article}
\usepackage[margin=1in]{geometry}
\usepackage{hyperref}
\usepackage{enumitem}
\usepackage{graphicx}
\usepackage{xcolor}

% -----------------------------------------------------
% bvraghav-tiet-ucs503p-article.sty
% -----------------------------------------------------

% Variable: Lab Instructor
\makeatletter \newcommand{\BvrSetLabInstructor}[1]{%
  \gdef\Bvr@LabInstructor{#1}}
% default
\newcommand{\Bvr@LabInstructor}{Dr. Raghav %
  B. Venkataramaiyer}
% public accessor
\newcommand{\BvrLabInstructorValue}{\Bvr@LabInstructor}
\makeatother

% Add subtitle
\usepackage{titling}
% -- Set title bold
\pretitle{\begin{center}\bfseries\fontsize{18bp}{18bp}\selectfont}
\makeatletter
\newcommand{\BvrSubtitle}[1]{%
  \posttitle{%
    \par\end{center}
    \begin{center}\large#1\end{center}
    \vskip 0.5em}%
}
\makeatother

% Hints in the section title(s)
\newcommand{\BvrHint}[1]{%
  {\normalfont\normalsize\itshape\hfill #1}}

% -----------------------------------------------------

\title{Deadstock Management System}

\BvrSetLabInstructor{Dr. Asif}

\BvrSubtitle{%
  Project Report for UCS503P  \\
  Submitted to: \BvrLabInstructorValue \\ \bfseries
  Thapar Institute of Engineering and Technology%
}

\author{
  Aryan Gupta, CSED%
  \thanks{Roll No: \texttt{1024030455}, %
    Email: \texttt{agupta\_be24\@thapar.edu}}
  \and
  Abhilakshya Puri, CSED%
  \thanks{Roll No: \texttt{1024030445}, %
    Email: \texttt{apuri\_be24\@thapar.edu}}
  \and
  Farhan Kansal, CSED%
  \thanks{Roll No: \texttt{1024030451}, %
    Email: \texttt{fkansal\_be24\@thapar.edu}}
}

\date{\today}

\begin{document}
\maketitle

\section{Higher-order goal%
  \BvrHint{\ldots secondary framing}}
This project aims to reduce capital wastage and material
waste in retail supply chains by giving unsold
(``deadstock'') inventory a structured path back into the
market. While the broader goal is sustainable, efficient
retail, the immediate engineering objective is a
deployable web platform that lets vendors track, verify,
and resell aging stock to other buyers.

\section{Time-to-value %
\BvrHint{\ldots primary heuristic}}
\begin{itemize}[leftmargin=0.7cm]
    \item We will deliver meaningful increments quickly by prototyping the core \emph{list-and-browse} workflow first, before layering on verification and compliance checks.
    \item We will prioritize fast feedback over perfection by using weekly iterations and CI-backed automation early to reduce release friction.
    \item Success is defined by what can be delivered and measured in the first iteration (a working vendor-to-buyer listing loop), then improved based on user feedback.
\end{itemize}

\section{Problem Statement}
An estimated 20--30\% of retail inventory value in India
sits idle as slow-moving or dead stock annually, and
over 2,00,000 licensed liquor and FMCG retailers alone
face recurring dead-stock write-offs each year. Vendors
currently have no structured way to recover this trapped
capital. Common pain points include:
\begin{itemize}[leftmargin=0.7cm]
    \item \textbf{No real-time visibility:} vendors track stock in spreadsheets or registers, so items are only noticed \emph{after} they've become deadstock.
    \item \textbf{No way to verify deadstock:} buyers cannot trust a listing without proof of condition, quantity, or authenticity, since none of this is independently confirmed today.
    \item \textbf{No resale marketplace:} even genuinely resellable stock has no dedicated channel connecting vendors with willing buyers.
    \item \textbf{Compliance risk:} regulated categories (e.g., licensed liquor) cannot resell freely; any resale channel must respect licensing rules.
\end{itemize}

\textbf{Why this matters now (contextual relevance):} unsold stock is routinely discarded, discounted below cost, or destroyed --- a direct loss to the vendor and to the environment --- while a trusted resale channel could recover that value instead.

\section{Proposed Solution %
\BvrHint{\ldots Software Proposition}}
\subsection{Overview}
Build a \textbf{web-based verified marketplace} where
vendors track, audit, and resell deadstock inventory to
other buyers, with the following components:
\begin{itemize}[leftmargin=0.7cm]
    \item \textbf{Vendor inventory dashboard:} digital stock logging with automated aging alerts flagging items drifting toward deadstock status.
    \item \textbf{Verification layer:} photo/batch proof plus rule checks confirming condition, quantity, and authenticity before a listing goes live.
    \item \textbf{Marketplace listings:} verified deadstock published, categorized, and priced for discovery by other buyers.
    \item \textbf{Transaction flow:} buyers browse trusted listings and purchase directly from vendors.
    \item \textbf{Compliance guardrails:} category-aware rules so regulated stock (e.g., licensed liquor) cannot be listed outside permitted channels.
\end{itemize}

\subsection{Core workflow}
\begin{enumerate}[leftmargin=0.7cm]
    \item Vendor logs stock digitally; the system flags items aging toward deadstock status.
    \item Vendor submits a deadstock item for verification (photo/batch proof).
    \item System applies rule checks (completeness, category compliance) and a human/automated review confirms the listing.
    \item Verified listing is published to the marketplace; buyers browse, filter, and purchase.
    \item Vendor and buyer view transaction and listing status through to completion.
\end{enumerate}

\subsection{Operational constraints for an educational, tier-2/3 setting}
\begin{itemize}[leftmargin=0.7cm]
    \item Keep the solution usable on low-to-mid range devices via responsive web design, since many vendors operate small retail outlets.
    \item Avoid dependence on advanced research algorithms (e.g., no computer-vision authenticity detection in scope); rely on structured proof plus rule-based checks.
    \item Design for deployability using common web hosting and managed databases.
\end{itemize}

\section{Solution Approach %
\BvrHint{\ldots Engineering focus}}
This is an engineering course project, so we focus on
scalable marketplace and workflow aspects rather than
novel algorithm research.

\subsection{Verification assistance %
\BvrHint{\ldots non-research}}
Verification will be checklist- and rule-based, not
ML-driven:
\begin{itemize}[leftmargin=0.7cm]
    \item Require photo/batch proof and structured fields (category, quantity, condition) at listing time.
    \item Apply deterministic rule checks (e.g., category-specific compliance rules, completeness checks) before publishing.
    \item Flag ambiguous or regulated-category listings for manual review; never auto-publish a regulated item without a pass.
\end{itemize}

\subsection{Web architecture %
\BvrHint{\ldots web-based solution}}
A typical 3-tier web architecture:
\begin{itemize}[leftmargin=0.7cm]
    \item \textbf{Frontend:} responsive UI for vendor dashboard, listing creation, and buyer marketplace/browse pages.
    \item \textbf{Backend API:} authentication, inventory CRUD, verification rule engine, listing and transaction logic.
    \item \textbf{Database:} vendors, buyers, inventory items, verification records, listings, transaction history.
\end{itemize}

\subsection{CI/CD and fast delivery enabler}
CI/CD will be used to:
\begin{itemize}[leftmargin=0.7cm]
    \item Automatically build and run tests on every change.
    \item Produce repeatable deployment artifacts to staging.
    \item Enable safe and frequent releases of incremental features.
\end{itemize}

\section{Evaluation Criterion %
\BvrHint{\ldots measurable and attributable}}
We define success using measurable metrics tied to the
core workflow.

\subsection{Primary evaluation metric}
\textbf{Time-to-listing (TTL):} time between a vendor
flagging an item as deadstock and it appearing as a
verified, live marketplace listing.
\begin{itemize}[leftmargin=0.7cm]
    \item Target: median TTL $\le$ 24 hours during pilot window.
    \item Attribution: measured from database timestamps (item flagged vs. listing published).
\end{itemize}

\subsection{Secondary evaluation metrics}
\begin{itemize}[leftmargin=0.7cm]
    \item \textbf{Recovery rate:} percentage of listed deadstock items sold within the pilot window.
    \item \textbf{Verification pass rate:} percentage of submitted listings passing rule checks without manual escalation.
    \item \textbf{Buyer trust:} buyer-reported confidence in listing accuracy using a simple 1--5 feedback question.
    \item \textbf{Reliability:} availability of staging/production for browsing and listing during business hours (e.g., $\ge$ 99\% uptime in pilot).
\end{itemize}

\subsection{Pilot validation plan %
\BvrHint{\ldots high-level}}
\begin{itemize}[leftmargin=0.7cm]
    \item Recruit 10--20 pilot vendors and a matching pool of buyers (simulated buyer roles if real ones are unavailable).
    \item Compare recovery rate and time-to-listing before/after within the iteration window, using baseline estimates from vendor interviews.
\end{itemize}

\section{Scalability %
\BvrHint{\ldots with foresight}}
The proposition should scale in both theoretical and
practical senses.

\begin{enumerate}[leftmargin=0.7cm]
    \item \textbf{Scalable (theoretically):} The system should support growth in number of vendors, listings, and concurrent buyers by:
    \begin{itemize}[leftmargin=0.7cm]
        \item decoupling components (frontend/backend),
        \item enabling pagination and indexing for common marketplace queries,
        \item using stateless API design.
    \end{itemize}

    \item \textbf{Operationally deployable (practically):} The system should run on common web hosting:
    \begin{itemize}[leftmargin=0.7cm]
        \item managed database,
        \item containerized or platform-hosted deployment,
        \item CI/CD pipeline for staging/production.
    \end{itemize}
\end{enumerate}

\section{Engine availability heuristic}
A mature set of ``engines'' (tooling/services) is
available to implement this as a web-based project:
\begin{itemize}[leftmargin=0.7cm]
    \item Web framework for REST APIs and templated/reactive pages.
    \item Managed relational database (for structured inventory/listing/transaction data).
    \item Authentication approach (token-based or session-based) for vendor and buyer accounts.
    \item Object storage for verification photos and batch-proof documents.
    \item Test framework for unit/integration tests.
    \item CI runner and staging deployment target.
\end{itemize}
We will use these mature components to focus on
marketplace design, verification correctness, and
measurement rather than inventing new infrastructure.

\section{Project scope and deliverables %
\BvrHint{\ldots iteration-friendly}}
\subsection{Initial deliverable %
\BvrHint{\ldots first iteration}}
\begin{itemize}[leftmargin=0.7cm]
    \item Vendor inventory entry + aging alert logic
    \item Listing creation with photo/proof upload
    \item Rule-based verification checklist (non-regulated categories)
    \item Buyer-facing browse/search page for live listings
    \item Basic logging and timestamps for TTL measurement
    \item CI: build + backend unit tests + linting
    \item CD to staging with a single command or pipeline job
\end{itemize}

\subsection{Subsequent deliverables}
\begin{itemize}[leftmargin=0.7cm]
    \item Regulated-category compliance rules and manual-review escalation path
    \item Transaction/purchase flow between buyer and vendor
    \item Vendor and admin dashboards for recovery-rate metrics
    \item Improved search/filtering and indexed queries for scale readiness
\end{itemize}

\section{Risks and mitigations}
\begin{itemize}[leftmargin=0.7cm]
    \item \textbf{Regulatory risk:} restrict regulated categories (e.g., licensed liquor) from auto-publishing; require manual compliance review before those listings go live.
    \item \textbf{Trust/fraud risk:} require photo/batch proof at listing time and keep an audit trail of verification decisions.
    \item \textbf{Vendor adoption:} keep listing creation to a minimal number of steps; provide clear aging alerts so vendors act before stock is a total write-off.
    \item \textbf{Evaluation measurement ambiguity:} instrument timestamps and define clear listing-status transitions before development begins.
\end{itemize}

\section{Summary}
This project proposes a deployable web marketplace that
turns unsold vendor stock into recovered revenue, by
combining inventory tracking, rule-based verification,
and a buyer-facing resale channel. It meets the course
criteria by emphasizing time-to-value through rapid
iterations, implementing CI/CD to support frequent safe
releases, and using measurable evaluation criteria
(especially time-to-listing and recovery rate). It
remains focused on engineering workflow, trust, and
scalability readiness rather than research-driven
novelty.

\end{document}